%--------------------------------------
% Caso de Uso: Añadir Empleado
%--------------------------------------

\begin{UseCase}{CU23}{Añadir Empleado}{
	Permitir que el Administrador registre un nuevo empleado en el sistema del Hotel para Perros.
}
	\UCitem{Versión}{\color{Gray}1.0}
	\UCitem{Autor}{\color{Gray}Vazquez Villagran Jorge}
	\UCitem{Supervisa}{\color{Gray}Beltran Saucedo Axel Alejandro}
	\UCitem{Actor}{Administrador}
	\UCitem{Propósito}{Registrar la información personal y laboral de un nuevo empleado.}
	\UCitem{Entradas}{Nombre, apellido paterno, apellido materno, email, teléfono, puesto, salario y fecha de contratación.}
	\UCitem{Origen}{Teclado}
	\UCitem{Salidas}{Empleado registrado y mensaje de confirmación.}
	\UCitem{Destino}{Pantalla de lista de empleados}
	\UCitem{Precondiciones}{
		El Administrador debe haber iniciado sesión.
	}
	\UCitem{Postcondiciones}{
		El nuevo empleado queda registrado en la base de datos con estado activo.
	}
	\UCitem{Errores}{
		Datos inválidos, email duplicado, permisos insuficientes.
	}
	\UCitem{Tipo}{Caso de uso primario}
	\UCitem{Observaciones}{
		Solo los usuarios con rol de administrador pueden añadir empleados.
	}
\end{UseCase}

%--------------------------------------
\begin{UCtrayectoria}
	\UCpaso[\UCactor] Selecciona la opción \IUbutton{Nuevo Empleado} en la Pantalla de Lista de Empleados.
	\UCpaso El sistema muestra la Pantalla de Registro de Empleado con el formulario vacío.
	\UCpaso[\UCactor] Captura la información del nuevo empleado.
	\UCpaso[\UCactor] Presiona el botón \IUbutton{Guardar}.
	\UCpaso El sistema valida que los datos ingresados sean correctos.
	\UCpaso El sistema verifica que el email no esté previamente registrado.
	\UCpaso El sistema registra al nuevo empleado en la base de datos.
	\UCpaso El sistema muestra el mensaje ``Empleado registrado correctamente''.
	\UCpaso El sistema regresa a la Pantalla de Lista de Empleados.
\end{UCtrayectoria}

%--------------------------------------
\begin{UCtrayectoriaA}{A}{Datos inválidos}
	\UCpaso El sistema muestra mensajes indicando los campos que contienen errores.
	\UCpaso Continúa en el paso 3 de la trayectoria principal.
\end{UCtrayectoriaA}

%--------------------------------------
\begin{UCtrayectoriaA}{B}{Email duplicado}
	\UCpaso El sistema muestra el mensaje ``El email ingresado ya está registrado''.
	\UCpaso Continúa en el paso 3 de la trayectoria principal.
\end{UCtrayectoriaA}

%--------------------------------------
\begin{UCtrayectoriaA}{C}{Permisos insuficientes}
	\UCpaso El sistema muestra el mensaje ``No cuenta con permisos para registrar empleados''.
	\UCpaso Termina el caso de uso.
\end{UCtrayectoriaA}

%--------------------------------------
\subsection{Puntos de extensión}
\UCExtenssionPoint{
	El Administrador desea asignar al empleado a servicios o citas.
}{
	Del paso 7 al paso 9.
}{
	Asignar Empleado a Servicio.
}

%--------------------------------------
% Fin del Caso de Uso
